\documentclass[11pt]{article}
\usepackage[margin=1in]{geometry}
\usepackage{amsmath,amssymb}
\usepackage[hidelinks]{hyperref}

\newcommand{\Lip}{\operatorname{Lip}_0}
\newcommand{\F}{\mathcal F}

\title{A Literature-Implied Answer to the Vector-Valued Lipschitz Octahedrality Question}
\author{arXiv automatic math run \texttt{fa\_banach\_001}}
\date{2026-06-18}

\begin{document}
\maketitle

\section*{Status}

\textbf{Status: literature-implied answer.}  This packet records that a
question from Langemets--Rueda Zoca, arXiv:1905.09061, is answered by combining
later published/preprint results.  It is not claimed as an original proof.

\section*{Source Question}

The source paper is J. Langemets and A. Rueda Zoca, \emph{Octahedral norms in
duals and biduals of Lipschitz-free spaces}, arXiv:1905.09061.  Problem 1.3 in
that paper asks:

\begin{quote}
Let \(M\) be a metric space such that \(\Lip(M)\) is octahedral and let \(X\)
be a non-trivial Banach space. Is then the space \(\Lip(M,X)\) octahedral?
\end{quote}

\section*{Answer}

Yes.  The answer is affirmative.

\section*{Identification}

The first ingredient is A. Rueda Zoca, \emph{Banach spaces which always produce
octahedral spaces of operators}, arXiv:2207.08717.  Example 4.3 states, using
arXiv:1905.09061, Proposition 3.3, that if \(\Lip(M)\) is octahedral, then
\(\Lip(M)\) is \(c_0\)-octahedral and consequently universally octahedral.
In particular,
\[
        L(Y,\Lip(M))
\]
is octahedral for every Banach space \(Y\).

The second ingredient is R. Medina and A. Rueda Zoca, \emph{A characterisation
of the Daugavet property in spaces of vector-valued Lipschitz functions},
arXiv:2305.05956.  Theorem 4.2 characterizes universal octahedral domains:
for a Banach space \(Z\), the following are equivalent:
\[
        L(Z,Y)\ \text{is octahedral for every Banach space }Y
\]
and
\[
        Z^*\ \text{is universally octahedral}.
\]
The theorem is stated in a slightly stronger form, allowing all subspaces of
\(L(Z,Y)\) that contain the finite-rank operators.

Apply Theorem 4.2 with \(Z=\F(M)\).  Since \(\F(M)^*=\Lip(M)\) and \(\Lip(M)\)
is universally octahedral by the preceding paragraph, \(L(\F(M),X)\) is
octahedral for every Banach space \(X\).  Finally, the standard linearization
isometry gives
\[
        \Lip(M,X) \cong L(\F(M),X).
\]
Therefore \(\Lip(M,X)\) is octahedral for every non-trivial Banach space
\(X\), answering Problem 1.3 affirmatively.

\section*{Scope and Limitations}

This packet only records the implication chain above.  It does not settle the
broader tensor-product question asking whether \(X\widehat{\otimes}_\pi Y\) is
octahedral whenever \(X\) is octahedral.

\section*{References}

\begin{itemize}
\item J. Langemets and A. Rueda Zoca, \emph{Octahedral norms in duals and
biduals of Lipschitz-free spaces}, arXiv:1905.09061.
\item A. Rueda Zoca, \emph{Banach spaces which always produce octahedral
spaces of operators}, arXiv:2207.08717.
\item R. Medina and A. Rueda Zoca, \emph{A characterisation of the Daugavet
property in spaces of vector-valued Lipschitz functions}, arXiv:2305.05956.
\end{itemize}

\end{document}
